\section{Evaluation}
\label{sec:experiment_results}
We separate participant observations of sign response and triggering from
current software verification of evidence handling. Historical repeated measures,
held-out speech tests, and deployment records characterize component behavior.

\paperinput{sections/validation-study}

\subsection{Historical Healthy-Participant Measurements}
Forty healthy participants completed five visits each, yielding 200 sessions
and 800 B/E/F/A records. Each module had 196 analyzable results, one positive
(0.51\%), and four insufficient outcomes (2.0\% of attempts).
Figure~\ref{fig:healthy} reports ICC(1,1) for one recorded output per module,
with participant-bootstrap intervals (5,000 resamples; seed 203).
B/E/F/A estimates were 0.140/0.266/0.784/0.679, indicating limited B/E
repeatability relative to between-participant variation.
The outputs were trunk-orientation change, binocular directional asymmetry,
mouth-corner delta, and arm-level difference, respectively.
These component outcomes do not estimate the current workflow's false-alert
probability. Balance baseline linkage and Face shoulder-width normalization
versus interocular normalization in the implementation require reconciliation.
The source of the historical E positive label (questionnaire or automatic
output) is unspecified; current automatic E classification is disabled. These
statistics describe historical repeatability, not current threshold reliability.

\begin{figure}[t]
\centering
\includegraphics[width=0.88\columnwidth]{healthy-repeatability.pdf}
\caption{Repeatability of one selected output per component, with 196 eligible measurements each. Bars show participant-bootstrap 95\% intervals for ICC(1,1).}
\label{fig:healthy}
\end{figure}
\subsection{Speaker-Independent Speech Evaluation}
The Mandarin Dysarthria Speech Corpus (MDSC) contains 18,630 recordings from 46 Mandarin speakers: 21 with dysarthria
and 25 controls \cite{gao2024mdsc}. Speaker-disjoint training, validation, and
test sets contained 33/7/6 speakers and 13,365/2,835/2,430 recordings.
The logistic model used Log-Mel, temporal, energy, and spectral descriptors. Its threshold was frozen at 0.735282 using validation
sensitivity of at least 0.90 and maximal specificity. Intervals resampled speakers.

Table~\ref{tab:model_performance} reports held-out results. Sensitivity fell
from 90.04\% in validation to 67.24\% in testing, while specificity rose from
81.36\% to 98.85\%. One speaker accounted for 313 of 398 misses
(Fig.~\ref{fig:speakers}). Only six independent speakers contributed to testing,
explaining the wide uncertainty despite 2,430 recordings. These results concern
prompted dysarthric speech in quiet conditions; household and acute-stroke
performance remain untested.

\begin{table}[htbp]
\centering\footnotesize
\caption{Held-out speech results: 2,430 recordings from six unseen speakers.
Intervals resample speakers, not individual recordings.}
\label{tab:model_performance}
\begin{tabular}{@{}lrr@{}}\toprule
Metric & Estimate & 95\% interval \\\midrule
Sensitivity & 67.24\% & 22.7--94.8\% \\
Specificity & 98.85\% & 98.0--99.8\% \\
Balanced accuracy & 83.05\% & 60.7--96.9\% \\
ROC--AUC & 0.9573 & --- \\
Brier score & 0.0983 & 0.022--0.227 \\
\bottomrule\end{tabular}
\end{table}
\begin{figure}[t]
\centering
\includegraphics[width=0.88\columnwidth]{mdsc-speaker-results.pdf}
\caption{Held-out speech performance at the frozen threshold, with 405 recordings per speaker. Control bars show specificity; dysarthric bars show sensitivity.}
\label{fig:speakers}
\end{figure}
\subsection{Pi--PC Deployment Pilot}
\label{sec:pipc_pilot}
Eight runs with two participants totaled 13~h, including one 8-h run,
with 99 events and 780 resource samples. The setup linked a 4-GB Raspberry
Pi~4, Camera V2, and ReSpeaker microphone HAT to an Intel NUC~13 Pro
with a Logitech C920. The pilot used the repository implementation.
Earlier and later runs showed Balance target counts of 4/7 and 6/7,
and Speech counts of 1/2 and 2/2. Negative-interval triggers were 2/11 and
1/11 for Balance, and 2/7 and 1/7 for Speech. Uncontrolled configuration
changes and two participants limit this comparison to a descriptive observation.
Supporting materials retain scenario counts.

During the 8-h run, CPU utilization averaged 43.1\%, RAM 1500~MB,
and temperature 69.4$^\circ$C (maximum 83.4$^\circ$C;
Fig.~\ref{fig:resources}). Records included one unavailability flag,
two data-loss flags, and one restart. Sampling coverage therefore does not
establish uninterrupted operation. Twenty logged \textit{soon} outcomes remain
unmapped to the current \textit{urgent} rule. Speech timing lacks capture-window
alignment, and the guided Balance record describes a questionnaire rather than
30-s standing acquisition. These records were not used to validate urgency
agreement or automatic-check latency. External alert delivery was not configured.

\begin{figure*}[t]
\centering
\includegraphics[width=0.72\textwidth]{resource-trace.pdf}
\caption{CPU and temperature sampled every 60~s during the 8-h run. Crosses mark offline records; dotted lines mark data loss; traces break at offline records.}
\label{fig:resources}
\end{figure*}
\subsection{Decision-Fusion and Interaction Verification}
An independent append-only reference agreed with the implementation across
1,800 F/A/S transitions involving measurements, retries, symptoms, and corrections;
previous report snapshots remained unchanged. In 15 constructed cases across
B/E/F/A/S, removing symptom retention, snapshot isolation, or correction history
broke the corresponding contract; the full implementation passed all cases.
These constructed removals isolate failure modes.

Four API regression scenarios additionally passed, with 24 requests recorded:
late audio after a symptom report, stale audio after reset, symptom correction
and persistent history readback, and B/E retry with unknown onset. They exercise
production API, session, reporting, and storage code using synthetic audio and
stub recognition, without participant sensing. Traces and source hashes accompany the supporting materials.

The reducer matched all $5^5=3,125$ specified evidence combinations.
Complete-first suppression differs in 2,070 cases by construction, not measured
benefit. Supporting materials retain the fixtures and protocols.
